% =========================================================================
% DOCUMENTCLASS
% =========================================================================
\documentclass[12pt,oneside]{book}

% =========================================================================
% METADATOS
% =========================================================================
\newcommand{\universidad}{Universidad de Buenos Aires (UBA)}
\newcommand{\carrera}{Notariado}
\newcommand{\materia}{Derecho Notarial y Registral}
\newcommand{\titulo}{Contratos Inmobiliarios}
\newcommand{\subtituloapunte}{Prevención del Lavado de Activos y Derechos Reales sobre Inmueble Ajeno}
\newcommand{\autor}{Isaac Edgar Camacho Ocampo}
\newcommand{\anio}{2024}

% =========================================================================
% PAQUETES
% =========================================================================
\usepackage[utf8]{inputenc}
\usepackage[T1]{fontenc}
\usepackage[spanish,es-nodecimaldot]{babel}

\usepackage[
    top=2.5cm,
    bottom=2.5cm,
    left=3cm,
    right=2.5cm,
    headheight=15pt
]{geometry}

\usepackage{amsmath}
\usepackage{amssymb}
\usepackage{mathtools}

\usepackage{graphicx}
\usepackage{float}
\usepackage{caption}
\usepackage{subcaption}

\usepackage{booktabs}
\usepackage{array}
\usepackage{longtable}
\usepackage{tabularx}

\usepackage{enumitem}

\usepackage{xcolor}
\usepackage[most]{tcolorbox}

\usepackage{fancyhdr}
\usepackage{ragged2e}

\graphicspath{
    {./img/}
    {./graficos/}
    {./figuras/}
}

\usepackage[
    colorlinks=true,
    linkcolor=blue!50!black,
    citecolor=green!40!black,
    urlcolor=blue!60!black,
    pdftitle={\titulo},
    pdfauthor={\autor},
    pdfsubject={Apuntes universitarios},
    bookmarks=true
]{hyperref}

% =========================================================================
% CONFIGURACIÓN
% =========================================================================
\setcounter{tocdepth}{2}
\setlength{\parindent}{0pt}
\setlength{\parskip}{0.5em}

\setlist[itemize]{
    topsep=0.3em,
    itemsep=0.2em
}
\setlist[enumerate]{
    topsep=0.3em,
    itemsep=0.2em
}

\clubpenalty=10000
\widowpenalty=10000

\pagestyle{fancy}
\fancyhf{}
\fancyhead[L]{\nouppercase{\leftmark}}
\fancyhead[R]{\materia}
\fancyfoot[C]{\thepage}
\renewcommand{\headrulewidth}{0.4pt}
\renewcommand{\footrulewidth}{0pt}

\fancypagestyle{plain}{
    \fancyhf{}
    \fancyfoot[C]{\thepage}
    \renewcommand{\headrulewidth}{0pt}
}

% =========================================================================
% COMANDOS
% =========================================================================
\newcommand{\abs}[1]{\left|#1\right|}
\newcommand{\norm}[1]{\left\lVert#1\right\rVert}
\newcommand{\vect}[1]{\mathbf{#1}}
\newcommand{\deriv}[2]{\frac{d#1}{d#2}}
\newcommand{\pderiv}[2]{\frac{\partial#1}{\partial#2}}

% =========================================================================
% ENTORNOS / CAJAS
% =========================================================================
\newtcolorbox{definition}[1]{
    enhanced,
    breakable,
    title={Definición: #1},
    fonttitle=\bfseries,
    colback=blue!3,
    colframe=blue!50!black,
    boxrule=0.8pt,
    arc=2mm
}

\newtcolorbox{important}{
    enhanced,
    breakable,
    title={Importante},
    fonttitle=\bfseries,
    colback=yellow!8,
    colframe=orange!70!black,
    boxrule=0.8pt,
    arc=2mm
}

\newtcolorbox{example}[1]{
    enhanced,
    breakable,
    title={Ejemplo: #1},
    fonttitle=\bfseries,
    colback=green!4,
    colframe=green!40!black,
    boxrule=0.8pt,
    arc=2mm
}

\newtcolorbox{formula}[1]{
    enhanced,
    breakable,
    title={Fórmula / Regla: #1},
    fonttitle=\bfseries,
    colback=purple!4,
    colframe=purple!50!black,
    boxrule=0.8pt,
    arc=2mm
}

\newtcolorbox{exercise}[1]{
    enhanced,
    breakable,
    title={Ejercicio: #1},
    fonttitle=\bfseries,
    colback=gray!5,
    colframe=gray!60!black,
    boxrule=0.8pt,
    arc=2mm
}

\newtcolorbox{note}{
    enhanced,
    breakable,
    title={Nota},
    fonttitle=\bfseries,
    colback=white,
    colframe=gray!50,
    boxrule=0.6pt,
    arc=2mm
}

% =========================================================================
% DOCUMENT
% =========================================================================
\begin{document}

% -------------------------------------------------------------------------
% PORTADA
% -------------------------------------------------------------------------
\begin{titlepage}
\thispagestyle{empty}

\begin{center}

\vspace*{1cm}

{\large \textsc{\universidad}}

\vspace{0.2cm}

{\large \carrera}

\vspace{3cm}

{\Huge \textbf{\titulo}}

\vspace{0.6cm}

{\LARGE \subtituloapunte}

\vspace{0.8cm}

{\large \materia}

\vspace{1cm}

\textit{Comprender los fundamentos de cada instituto jurídico es la base para ejercer la función notarial con criterio y responsabilidad.}

\vspace{1.2cm}

\rule{0.70\textwidth}{0.5pt}

\vspace{0.5cm}

{\large Apuntes de estudio}

\vspace{0.5cm}

\rule{0.70\textwidth}{0.5pt}

\vfill

{\large \autor}

\vspace{0.3cm}

{\large \anio}

\end{center}

\vfill

\begin{flushleft}
\textit{Este es un modesto aporte a los estudiantes de ``Derecho Notarial y Registral'', no pretende sustituir el material oficial de la materia.}
\end{flushleft}

\end{titlepage}

% -------------------------------------------------------------------------
% ÍNDICE
% -------------------------------------------------------------------------
\tableofcontents

% =========================================================================
% CONTENIDO
% =========================================================================

\chapter{Prevención del lavado de activos ante la Unidad de Información Financiera (UIF)}

\section{El escribano como sujeto obligado}

El escribano reviste, en el ejercicio de la función notarial, el carácter de \textbf{sujeto obligado} ante la Unidad de Información Financiera (UIF). Esta condición impone un deber especial de cuidado en la redacción de los instrumentos públicos, particularmente en lo relativo a las manifestaciones sobre la entrega de dinero entre las partes.

La inclusión de cláusulas del tipo \textit{``en este acto ante mí, las partes entregan la suma de\dots''} compromete al escribano como testigo directo del traspaso de fondos. Por ese motivo, se recomienda evitar dicha fórmula cuando no exista certeza sobre la trazabilidad del dinero involucrado en la operación.

\begin{important}
El escribano debe evitar presenciar directamente la entrega del dinero en una compraventa. Cuando el inmueble no permite retirarse del acto de entrega, una alternativa práctica es dirigir la atención hacia otro lugar durante ese momento. Esta conducta no constituye únicamente una práctica ética: es una precaución frente a eventuales situaciones de lavado de activos.
\end{important}

\begin{note}
La sola circunstancia de no estar presente en el momento material de la entrega del dinero no exime al escribano de responsabilidad, ya que la obligación de control sobre el origen de los fondos subsiste independientemente de la fórmula empleada en el acto.
\end{note}

\section{Determinación del origen de los fondos: un caso de trazabilidad incompleta}

\begin{example}{Escritura impugnada por origen de fondos}
Una persona con residencia en España realizó una compraventa inmobiliaria en Argentina. Con anterioridad, había retirado del país una suma de USD 50.000 a través del aeropuerto de Ezeiza. Al momento de justificar el origen de esos fondos, presentó la escritura de venta de la casa de su madre; sin embargo, ese instrumento constaba de una suma en pesos, sin justificativo alguno de la conversión de esos pesos a dólares. La operación quedó bajo sospecha porque la cadena de trazabilidad del dinero se encontraba incompleta.
\end{example}

De este caso se desprende que la ausencia de una cláusula que documente expresamente la entrega de dinero no libera por completo al escribano de responsabilidad: si la operación resulta sospechosa por falta de trazabilidad, corresponde directamente no realizarla.

\section{Criterios para determinar el carácter sospechoso de una operación}

\subsection{El criterio de sospecha no se basa en el origen nacional del cliente}

La nacionalidad del comprador o vendedor no constituye, por sí sola, un indicador válido de sospecha. Existen otros elementos objetivos que deben tomarse en cuenta para calificar una operación como sospechosa, y la sola procedencia extranjera del cliente no resulta suficiente para ese juicio.

\subsection{Indicadores reales de sospecha}

\begin{table}[H]
\centering
\caption{Indicadores de sospecha en una operación inmobiliaria}
\begin{tabularx}{\textwidth}{>{\raggedright\arraybackslash}X >{\raggedright\arraybackslash}X c}
\toprule
\textbf{Indicador} & \textbf{Ejemplo práctico} & \textbf{¿Sospechoso?} \\
\midrule
Capacidad contributiva & Persona sin actividad registrada que compra un inmueble de alto valor & Sí \\
Nacionalidad solamente & Comprador ruso, americano, uruguayo & No \\
Aspecto físico / vestimenta & Persona con vestimenta asociada a un origen determinado & No \\
Capacidad contributiva acreditada & Persona con varios negocios registrados que paga IVA y Ganancias & No \\
Exhibición de riqueza & Persona con bienes de lujo, sin justificación de fondos & Mayor sospecha \\
Países en lista negra & Irán, Myanmar, Corea del Norte & Sí -- zona negra UIF \\
Países en lista gris & Países que entraron y salieron de listas del GAFI & Mayor recaudo \\
\bottomrule
\end{tabularx}
\end{table}

\begin{example}{Prejuicio por apariencia frente a análisis objetivo}
Una persona vestida de manera diferente a la habitual en una inmobiliaria puede generar desconfianza injustificada en el personal del lugar; sin embargo, la apariencia no constituye motivo válido de sospecha por sí misma. Si esa misma persona pretendiera adquirir un inmueble por una suma elevada, correspondería entonces analizar la trazabilidad del origen de los fondos y, eventualmente, solicitar un certificado de origen, con independencia de su vestimenta, ya que podría tratarse de alguien con varios negocios registrados que tributa regularmente. En sentido inverso, una persona que exhibe bienes de lujo sin justificación de ingresos amerita un control más riguroso que aquella cuya apariencia no resulta llamativa.
\end{example}

\section{Criptomonedas en operaciones inmobiliarias}

\subsection{El Bitcoin como origen de fondos}

Cuando el dinero utilizado para una compra proviene de la venta previa de criptomonedas, se presenta la dificultad de acreditar dicho origen ante el escribano. Si los bitcoins se encuentran en una \textit{billetera fría} (dispositivo físico sin conexión permanente a internet), esa circunstancia no hace, por sí sola, sospechosa a la operación. El problema principal no radica en el medio utilizado, sino en la dificultad de determinar el valor exacto de la criptomoneda al momento de otorgarse la escritura.

\subsection{La volatilidad y la necesidad de una moneda estable}

Para determinar el pago del impuesto de sellos y el valor de la venta es necesario fijar un monto en un momento determinado. Dado que el valor de una criptomoneda puede variar entre el inicio de la lectura de la escritura y el momento de su firma, resulta necesaria una moneda estable como referencia para la operación.

\subsection{Calificación jurídica del pago en criptomonedas}

Cuando el pago de una operación inmobiliaria se realiza mediante criptomonedas, corresponde distinguir con precisión la figura jurídica aplicable, evitando encuadrar el acto en un tipo contractual que no le corresponde.

\begin{table}[H]
\centering
\caption{Figuras jurídicas según el medio de pago}
\begin{tabularx}{\textwidth}{>{\raggedright\arraybackslash}X >{\raggedright\arraybackslash}X >{\raggedright\arraybackslash}X}
\toprule
\textbf{Figura jurídica} & \textbf{Descripción} & \textbf{Moneda} \\
\midrule
Compraventa & Transferencia a título oneroso con precio en dinero & Moneda de curso legal \\
Permuta & Intercambio de una cosa por otra cosa & No hay precio en dinero \\
Dación en pago & Se extingue una deuda entregando una cosa distinta & Equivalente en bienes \\
Pago en especie & Se paga con bienes no dinerarios (incluye criptoactivos) & Bitcoin, USD, etc. \\
\bottomrule
\end{tabularx}
\end{table}

\begin{important}
El pago realizado con criptomonedas no debe calificarse como compraventa en sentido estricto, en tanto la moneda de curso legal en la República es la única que reviste ese carácter para dicha figura. El pago con bienes no dinerarios corresponde encuadrarse como dación en pago. Forzar el encuadre de un acto en un tipo contractual o societario que no le corresponde constituye un error técnico que debe evitarse.
\end{important}

\section{La importancia del encuadre jurídico correcto}

\subsection{Caso: gestión de negocios y sociedad inexistente}

\begin{example}{La sociedad no constituida al momento de la compra}
Una persona adquirió un inmueble en carácter de gestor de negocios para una sociedad que, oportunamente, aceptaría esa compra. El problema fue que dicha sociedad no había sido siquiera constituida al momento de la adquisición; estaba conformada por dos cuñadas. Falleció el marido de una de ellas sin que la sociedad hubiera aceptado la compra, por lo que el inmueble ingresó en la sucesión. La viuda, con cuatro hijos menores, solicitó autorización judicial para vender. El mismo juez que autorizó la venta luego dispuso que no se inscribiera la escritura. El caso llegó a Cámara, que revocó la decisión de primera instancia.
\end{example}

El encuadre incorrecto de un contrato —adquirir en carácter de gestor de negocios de una sociedad que en ese momento no existía ni disponía de los fondos para adquirir el bien— generó un conflicto que se prolongó en distintas instancias judiciales.

\subsection{El título del acto frente a su contenido}

\begin{important}
Lo determinante en una escritura es su contenido —esto es, que incluya todos los elementos que le corresponden según la figura jurídica de que se trate— y no el título o rótulo que se le asigne. Sin embargo, rotular incorrectamente un acto puede generar interpretaciones erróneas: si se denomina ``cesión'' a lo que en sustancia constituye una donación —por tratarse de una transmisión a título gratuito—, corresponde igualmente interpretarla como donación, con las consecuencias jurídicas propias de esta figura.
\end{important}

\subsection{Contratos informales, instrumentos en blanco y comunicaciones tácitas}

\begin{note}
El escribano no debería certificar una firma sobre un instrumento en blanco. Asimismo, corresponde tener presente que pueden existir contratos sin aceptación expresa —contratos tácitos— derivados incluso de comunicaciones informales, como una respuesta afirmativa enviada por mensajería instantánea.
\end{note}

\section{Operaciones que deben reportarse a la UIF}

\begin{important}
\textbf{Umbral de reporte UIF (artículo referenciado en clase):} deben reportarse las compraventas inmobiliarias superiores a 750 salarios mínimos vitales y móviles (SMVyM), las cesiones de participaciones societarias y las participaciones en fideicomisos.
% TODO: contenido incompleto o ambiguo en las notas originales — no se identificó en el apunte el número preciso de la norma que fija el umbral de 750 SMVyM.
\end{important}

\subsection{Actividades específicas frente a actividades especiales}

Las actividades especiales deben tramitarse mediante protocolo, mientras que los contratos celebrados en forma privada no se reportan ni se consideran actividades específicas. Esta distinción resulta relevante para determinar el alcance de la obligación de informar.

\begin{table}[H]
\centering
\caption{Actos notariales y su deber de información ante la UIF}
\begin{tabularx}{\textwidth}{>{\raggedright\arraybackslash}X c >{\raggedright\arraybackslash}X}
\toprule
\textbf{Actividad} & \textbf{¿Informa a la UIF?} & \textbf{Razón} \\
\midrule
Compraventa inmobiliaria por escritura ($\geq$ 750 SMVyM) & Sí & Actividad específica con intervención notarial \\
Boleto de compraventa por instrumento privado & No & No es acto notarial \\
Boleto de compraventa por escritura pública & Sí & Es un acto protocolar notarial \\
Constitución de sociedad (escritura) & Sí & Actividad específica notarial \\
Aumento de capital (protocolización de asamblea) & No & El acto lo tomó la asamblea, no el escribano \\
Cesión de cuotas / venta de acciones & Sí & Hay participación real del escribano en la transferencia \\
Cambio de directorio (art.\ 60, Ley 19.550) & No & Es una inscripción registral, no una transmisión \\
Modificación de estatuto & No & No implica transferencia de fondos \\
\bottomrule
\end{tabularx}
\end{table}

\subsection{El aumento de capital y la cesión de cuotas}

El aumento de capital constituye una decisión asamblearia que no pasa por la intervención directa del escribano; lo máximo que este realiza es la protocolización de dicho acto, el cual en sí mismo no reviste carácter notarial sino societario. Situación distinta se presenta cuando, a partir de ese aumento de capital, se produce una cesión de cuotas: en tal caso corresponde informar la operación, en tanto resulta necesario conocer el origen del dinero involucrado.

\section{Tipos de dinero y obligaciones del escribano}

La profesión notarial constituye una actividad de alto riesgo, en la que el escribano responde con su patrimonio personal por el incumplimiento de sus obligaciones. Resulta necesario distinguir tres categorías de dinero:

\begin{table}[H]
\centering
\caption{Tipos de dinero y conducta del escribano}
\begin{tabularx}{\textwidth}{>{\raggedright\arraybackslash}X >{\raggedright\arraybackslash}X >{\raggedright\arraybackslash}X}
\toprule
\textbf{Tipo de dinero} & \textbf{Descripción} & \textbf{Conducta del escribano} \\
\midrule
Dinero blanco & De origen completamente legal y declarado & Procede sin restricciones \\
Dinero negro & No declarado fiscalmente, pero no necesariamente ilícito & Evalúa el riesgo; puede no realizar la operación \\
Dinero sucio & Proveniente de actividades ilícitas (lavado, narcotráfico, etc.) & No la realiza; denuncia obligatoria \\
\bottomrule
\end{tabularx}
\end{table}

El dinero negro no equivale necesariamente a dinero sucio, por lo que ambas categorías no deben tratarse de igual manera. Ante un caso de dinero sucio, corresponde directamente no realizar la operación.

\begin{important}
El escribano cuenta con 24 horas para denunciar ante la UIF una operación sospechosa detectada. La omisión del reporte puede acarrear sanciones penales y la pérdida de la habilitación profesional.
\end{important}

\begin{note}
El deber de conocer al cliente no equivale a una obligación de investigación exhaustiva sobre cada persona con la que se opera. El escribano no reviste el carácter de investigador; su deber se circunscribe al análisis razonable de los elementos objetivos disponibles al momento de la operación.
\end{note}

\section{Cuadro Cornell: prevención del lavado de activos}

\begin{table}[H]
\centering
\begin{tabularx}{\textwidth}{
    >{\raggedright\arraybackslash}p{0.30\textwidth}
    >{\raggedright\arraybackslash}X
}
\toprule
\textbf{Preguntas / palabras clave} & \textbf{Notas / respuestas} \\
\midrule

\textbf{¿Qué es la UIF y por qué obliga al escribano?} &
La Unidad de Información Financiera (UIF) es el organismo argentino de control anti-lavado. El escribano es un sujeto obligado legalmente: debe identificar al cliente, controlar el origen de los fondos y reportar las operaciones que superen ciertos umbrales o resulten sospechosas. \\

\textbf{¿Qué diferencia hay entre dinero blanco, negro y sucio?} &
Blanco: declarado y de origen lícito. Negro: no declarado fiscalmente, pero no necesariamente proveniente de un delito. Sucio: de origen ilícito. El escribano trabaja sin restricciones con el blanco, evalúa caso por caso el negro y debe rechazar y denunciar el sucio. \\

\textbf{¿Cuándo se reporta una operación a la UIF?} &
Se reportan las compraventas inmobiliarias por escritura que superen 750 SMVyM, las cesiones de participaciones societarias y los fideicomisos. No se reportan los boletos privados, los aumentos de capital por asamblea (solo protocolizados), los cambios de directorio ni las modificaciones de estatuto. \\

\textbf{¿Qué hace sospechosa a una operación?} &
No la nacionalidad ni la apariencia del cliente. Sí: la falta de capacidad contributiva para la operación, la ausencia de trazabilidad del dinero, la pertenencia del cliente a listas negras o grises de la UIF, o la incongruencia entre la operación y los activos declarados. \\

\textbf{¿Qué es el artículo 60 de la Ley 19.550 en materia societaria?} &
Es la norma que regula la inscripción de cambios de directorio, administradores y otros actos societarios internos en el Registro Público. El escribano puede intervenir protocolizando el acta, pero el acto en sí es de la sociedad y no implica reporte a la UIF, porque no hay transmisión de fondos ni de participaciones. \\

\midrule

\multicolumn{2}{p{\textwidth}}{
\textbf{Resumen:} El escribano actúa como guardián del sistema financiero legal: debe conocer a su cliente, evaluar la trazabilidad de los fondos y reportar ante la UIF las operaciones que superen los umbrales establecidos o presenten indicadores objetivos de sospecha. Esta obligación no puede cumplirse desde el prejuicio —la nacionalidad o la apariencia del cliente no son, por sí mismas, indicadores de sospecha— sino desde el análisis técnico de la capacidad contributiva, la coherencia de la operación y la verificación del origen del dinero.
} \\

\bottomrule
\end{tabularx}
\end{table}

% =========================================================================
\chapter{Derechos reales sobre inmueble ajeno}

\section{Usufructo: facultades, transmisión y estrategias familiares}

\subsection{Donación con reserva de usufructo}

\begin{example}{El padre que dona la nuda propiedad y se reserva el usufructo}
Un padre transmite la nuda propiedad del inmueble a sus hijos, pero desea continuar viviendo en él. Se reserva el usufructo, de modo que la nuda propiedad pasa a los hijos mientras el padre puede seguir percibiendo los frutos del bien.
\end{example}

\begin{important}
Consecuencia práctica: si el donante no se reserva el usufructo y el donatario fallece antes que él, el inmueble ingresa en la sucesión del donatario, y el donante puede quedar sin ningún derecho sobre el bien.
\end{important}

\subsection{Transmisión onerosa del usufructo}

Cuando el titular de un usufructo gratuito desea que los nudos propietarios vendan el inmueble sin renunciar a su participación económica, existe la posibilidad de transmitir en forma onerosa el propio usufructo de manera simultánea a la venta de la nuda propiedad.

\begin{example}{Venta simultánea de la nuda propiedad y del usufructo}
Un padre de 70 años se había reservado el usufructo gratuito de un inmueble al ceder a sus hijas la nuda propiedad. Ante la posibilidad de que ellas quisieran vender el bien, la estrategia consiste en que las hijas vendan la nuda propiedad y el padre transmita en forma onerosa su usufructo en el mismo acto. De esta manera, el padre recibe una parte del precio y las hijas reciben la otra parte.

Alternativamente, el padre puede transferir en forma onerosa el usufructo del inmueble original y adquirir, también en forma onerosa, el usufructo sobre un nuevo inmueble adquirido por las hijas, de modo de continuar contando con un lugar donde vivir mientras estas disponen del inmueble anterior.
\end{example}

\section{Derecho de uso: alcance y limitaciones}

\subsection{Características del derecho de uso frente al usufructo}

\begin{table}[H]
\centering
\caption{Derecho de uso y usufructo: comparación}
\begin{tabularx}{\textwidth}{>{\raggedright\arraybackslash}X >{\raggedright\arraybackslash}X >{\raggedright\arraybackslash}X}
\toprule
\textbf{Aspecto} & \textbf{Derecho de uso} & \textbf{Usufructo} \\
\midrule
Destino del inmueble & Solo para el uso personal del titular y su familia & Puede alquilarlo y percibir frutos \\
Transferibilidad & No transferible & Transferible onerosamente \\
Robustez & Menor & Mayor \\
Personas que pueden habitar & El titular y su núcleo familiar & El titular y quien él decida \\
\bottomrule
\end{tabularx}
\end{table}

\begin{example}{Derecho de uso para uso familiar estacional}
Una persona posee una vivienda en Mar del Plata que utiliza únicamente en temporada. Su hermano trabaja en esa ciudad durante el año, por lo que le otorga un derecho de uso para que pueda instalarse en esos períodos. El titular del derecho puede usar la vivienda para necesidades de su familia durante ese tiempo, pero no puede alquilarla ni disponer de ella.
\end{example}

\subsection{Límites prácticos del derecho de uso: el caso de la cochera}

Cuando se pretende constituir un derecho de uso sobre una unidad complementaria, como una cochera, corresponde verificar si esa unidad admite el ejercicio de un derecho real sobre una parte indivisa cuando el beneficiario no reviste el carácter de titular del edificio.

\begin{note}
Para situaciones de uso temporal y acotadas en el tiempo —por ejemplo, cuatro meses al año— no resulta conveniente constituir e inscribir un derecho real de uso en el registro. Una figura más adecuada es el contrato de comodato, que resulta más flexible, no requiere inscripción registral y permite recuperar el uso del inmueble si el comodante cambia de opinión.
\end{note}

\section{Derecho de habitación: protección de la vivienda personal}

\subsection{Características del derecho de habitación}

\begin{table}[H]
\centering
\caption{Derecho de habitación, derecho de uso y usufructo}
\begin{tabularx}{\textwidth}{>{\raggedright\arraybackslash}X c c c}
\toprule
\textbf{Aspecto} & \textbf{Habitación} & \textbf{Uso} & \textbf{Usufructo} \\
\midrule
¿Puede alquilar? & No & No (solo para sí) & Sí \\
¿Puede vivir con pareja? & No & Sí (familia) & Sí \\
¿Es transferible? & No & No & Sí \\
Robustez & Menor & Media & Mayor \\
\bottomrule
\end{tabularx}
\end{table}

\textbf{Uso típico:} el derecho de habitación se orienta a asegurar techo a una persona vulnerable; el derecho de uso, al uso personal con familia; y el usufructo, a la explotación económica del bien junto con su goce como vivienda.

\begin{example}{Protección de quien asistió a una persona mayor}
Una persona mayor había sido cuidada por otra. La hija de la primera no deseaba que quien había asistido a su madre quedara sin hogar, por lo que le constituyó un derecho de habitación. Bajo esta figura, la beneficiaria solo puede habitar el inmueble ella sola, sin sumar a otros integrantes.
\end{example}

\begin{important}
Corresponde evaluar con precisión cuál es el objetivo perseguido antes de constituir un usufructo cuando en realidad se pretende otorgar un derecho más limitado, como el de habitación, ya que el usufructo concede facultades considerablemente mayores a las que en ocasiones se desea otorgar.
\end{important}

\begin{example}{Alojamiento de un alumno de intercambio}
El titular de un departamento que no utiliza lo destina a un alumno que llega a estudiar durante un año desde el extranjero. Se le constituye un derecho de habitación para que resida en el inmueble durante ese período; bajo esta figura solo puede vivir el beneficiario, sin poder traer a otras personas ni alquilar el bien. A diferencia del comodato, donde el comodatario podría alojar a un tercero, el derecho de habitación es personalísimo, por lo que resulta la figura más precisa cuando el objetivo es que habite exclusivamente el beneficiario.
\end{example}

\section{Derecho de superficie: construcción sobre terreno ajeno}

\begin{important}
El derecho de superficie es el más robusto de los derechos reales sobre cosa ajena. Permite construir o plantar sobre un inmueble ajeno durante un plazo máximo fijado por el Código Civil y Comercial. Al vencer el plazo, lo construido revierte al propietario del terreno sin indemnización, salvo pacto en contrario.
\end{important}

\begin{example}{Salón de fiestas sobre terreno ajeno}
Una constructora necesita un predio para levantar un salón de fiestas. Se contacta con el titular de un terreno, quien le otorga el derecho de superficie por un plazo de 40 a 50 años para construir y explotar el salón; al finalizar el plazo, lo construido vuelve al titular del terreno. Una ventaja adicional del derecho de superficie frente al usufructo es que el superficiario puede, a su vez, alquilar o subarrendar lo construido.
\end{example}

\begin{example}{Construcción de departamentos por una empresa sobre terreno ajeno}
Una empresa desea construir departamentos sobre un terreno que no le pertenece. El propietario no está dispuesto a vender el terreno, pero acepta que la empresa construya y sea titular de lo edificado durante 50 años. Se constituye entonces un derecho de superficie, que permite construir sobre suelo ajeno; al vencer el plazo, si no se renueva, lo construido vuelve al titular del terreno.
\end{example}

\section{Comparación general de los derechos reales sobre cosa ajena}

\begin{table}[H]
\centering
\caption{Cuadro comparativo de derechos reales sobre inmueble ajeno}
\begin{tabularx}{\textwidth}{>{\raggedright\arraybackslash}X >{\raggedright\arraybackslash}X >{\raggedright\arraybackslash}X >{\raggedright\arraybackslash}X}
\toprule
\textbf{Derecho real} & \textbf{Principales facultades} & \textbf{Limitaciones} & \textbf{Uso típico} \\
\midrule
Usufructo & Usar, gozar, percibir frutos, ceder o alquilar & No puede deteriorar ni alterar el destino & Donación con reserva de renta; protección del donante \\
Uso & Usar el inmueble para sí y su familia & No puede arrendarlo; no transferible & Uso familiar temporal o estacional \\
Habitación & Habitar personalmente el inmueble & No puede albergar a terceros distintos de sí mismo & Protección de vivienda de persona vulnerable \\
Superficie & Construir, plantar y explotar sobre suelo ajeno & Plazo máximo legal; al vencer, revierte & Proyectos de construcción o explotación a largo plazo \\
\bottomrule
\end{tabularx}
\end{table}

\section{Usufructo vitalicio, legítima y colación: análisis de un fallo}

\subsection{El caso analizado}

\begin{example}{Usufructo vitalicio consentido por los hermanos}
Una mujer con cuatro hijos constituyó a favor del más pequeño un usufructo vitalicio sobre una propiedad, reservándose ella la nuda propiedad. Todos los hermanos firmaron consintiendo dicho acto. Transcurridos cuatro años, la mujer falleció. Con posterioridad, los hermanos pretendieron dejar sin efecto el usufructo del hermano menor, invocando el derecho de colación y la afectación de la porción legítima. Cabe agregar que ese hermano menor quedó luego incapacitado y actúa mediante representante.
\end{example}

El fallo de primera instancia, confirmado por la Cámara, rechazó el pedido de los hermanos por dos razones:

\begin{enumerate}
\item \textbf{Teoría de los actos propios:} los hermanos habían consentido originalmente el acto, por lo que no podían impugnarlo con posterioridad.
\item \textbf{Validez del usufructo vitalicio:} el usufructo había sido otorgado en forma vitalicia y válida, por lo que no podía revocarse unilateralmente; la nuda propiedad quedó en cabeza de los hermanos.
\end{enumerate}

\subsection{El usufructo frente a la porción legítima}

\begin{note}
% TODO: contenido incompleto o ambiguo en las notas originales — el apunte recoge, en este punto, una valoración personal del docente sobre la institución de la legítima (manifiesta estar en contra de su existencia), la cual se consigna aquí como opinión expresada en clase y no como contenido normativo vigente.
El docente señala que, si bien esa fue su postura personal expresada en clase, el caso ilustra una manera en que puede vulnerarse una norma de orden público inserta en el Código: la protección de la porción legítima de los herederos forzosos.
\end{note}

\begin{example}{Usufructo constituido a favor de cónyuge joven}
Un hombre de 90 años, con hijos de 60 años, contrae matrimonio con una mujer de 30 años y le constituye el usufructo sobre su patrimonio. Los hijos conservan la nuda propiedad, pero resulta probable que la cónyuge los sobreviva por un extenso período, de modo que los hijos quedarían con una nuda propiedad vacía de contenido económico durante ese lapso.
\end{example}

\begin{important}
El Código Civil y Comercial no distingue en función de la edad de las partes ni de probabilidades de supervivencia al regular la constitución del usufructo. Sin perjuicio de ello, cuando se constituye un usufructo sobre la totalidad del patrimonio a favor de una persona con expectativa de vida considerablemente mayor a la de los herederos forzosos, puede estarse vulnerando en los hechos la protección de la porción legítima, en tanto los herederos quedan con una nuda propiedad inaprovechable durante un período prolongado.
\end{important}

\section{Comodato frente a derechos reales: una distinción práctica}

Encasillar una situación dentro de un derecho real o de un tipo contractual determinado cuando ello no resulta necesario puede ser perjudicial para el cliente. Por ejemplo, para un uso acotado a pocos meses no resulta conveniente constituir e inscribir un derecho real de uso en el registro.

\begin{table}[H]
\centering
\caption{Comodato, derecho de uso y locación temporal}
\begin{tabularx}{\textwidth}{>{\raggedright\arraybackslash}X >{\raggedright\arraybackslash}X c >{\raggedright\arraybackslash}X >{\raggedright\arraybackslash}X}
\toprule
\textbf{Figura} & \textbf{Tipo} & \textbf{Inscripción} & \textbf{Duración típica} & \textbf{Recomendado para} \\
\midrule
Derecho de uso & Derecho real & Sí (registro) & Mediano/largo plazo & Uso permanente o estable \\
Comodato & Contrato personal & No & Corto/mediano plazo & Préstamo temporal sin costo \\
Locación temporal & Contrato personal & No (salvo más de 3 años) & Meses o años & Uso oneroso con plazo determinado \\
\bottomrule
\end{tabularx}
\end{table}

\begin{note}
El comodato es un contrato de préstamo de uso, distinto de los derechos reales antes analizados. Distinguir con precisión cuándo corresponde constituir un derecho real y cuándo basta un vínculo contractual personal resulta central para un asesoramiento notarial adecuado.
\end{note}

\section{Cuadro Cornell: derechos reales sobre inmueble ajeno}

\begin{table}[H]
\centering
\begin{tabularx}{\textwidth}{
    >{\raggedright\arraybackslash}p{0.30\textwidth}
    >{\raggedright\arraybackslash}X
}
\toprule
\textbf{Preguntas / palabras clave} & \textbf{Notas / respuestas} \\
\midrule

\textbf{¿Qué figuras pueden usarse para garantizar el uso de un inmueble sin transferirlo?} &
El usufructo (el más amplio: uso, goce y percepción de frutos; transferible), el derecho de uso (solo para el titular y su familia; no transferible), el derecho de habitación (solo para que viva el titular; no transferible, personalísimo) y el derecho de superficie (construir sobre suelo ajeno por plazo fijo; el más robusto de todos). \\

\textbf{¿Cuándo el donante debe reservar el usufructo?} &
Siempre que quiera seguir usando el inmueble o percibiendo sus frutos tras la donación. Si dona sin reservarlo y el donatario fallece, el inmueble pasa a los herederos de este y el donante queda sin derechos. La reserva de usufructo le asegura el uso vitalicio o por el plazo determinado. \\

\textbf{¿Puede el padre transmitir onerosamente su usufructo para que sus hijos vendan el inmueble?} &
Sí. El usufructo puede transmitirse onerosamente. El mecanismo consiste en que las hijas vendan la nuda propiedad y el padre transmita simultáneamente su usufructo. El comprador adquiere el pleno dominio y el padre recibe una parte del precio, manteniendo así su participación económica en la operación. \\

\textbf{¿Cuándo conviene usar el comodato en lugar de un derecho real?} &
Cuando la situación es temporal, informal y flexible. El comodato no requiere escritura ni inscripción registral, se adapta mediante el contrato y no genera un derecho real que luego deba cancelarse formalmente. Es la alternativa práctica cuando el uso del inmueble se extiende por meses o responde a situaciones ocasionales. \\

\textbf{¿Puede un usufructo vulnerar la porción legítima de los herederos?} &
Depende. Si el donante o testador constituye usufructos sobre la totalidad del patrimonio a favor de personas con expectativa de sobrevivir a los herederos forzosos, puede generarse una lesión a la porción legítima. El Código no distingue por edades ni probabilidades, pero en la práctica se presentan casos donde los herederos quedan con una nuda propiedad inaprovechable durante décadas. \\

\midrule

\multicolumn{2}{p{\textwidth}}{
\textbf{Resumen:} El ejercicio notarial exige dominar el instrumental de los derechos reales sobre cosa ajena como herramienta de asesoramiento. El usufructo, el uso, la habitación y la superficie constituyen soluciones concretas a necesidades reales de las personas. Elegir el instrumento equivocado puede privar a un donante de su hogar, otorgar a un tercero más derechos de los que se quisieron conceder, o dejar a los herederos con una nuda propiedad inútil durante décadas. Comprender la lógica de cada figura permite asesorar con precisión, proteger a las partes y hacer prevalecer la voluntad real de los clientes por encima de fórmulas predeterminadas.
} \\

\bottomrule
\end{tabularx}
\end{table}

\end{document}
